
\subsection{Contributions}
\paragraph{The theoretical importance of culture-talk.}
This paper contributes to the study of culture and stratification in several ways.  Analytically, I highlight the theoretical and empirical importance of the phenomenon that I have referred to as \emph{culture-based talk}. I began with the Simmelian legacy on the importance of seemingly trivial conversational content for the emergence of ``sociability'' (or the ``play form'' of association) and connected it to the more contemporary stress by cultural stratification theorists on ``entertainment talk'' as the cement that holds status groups together as micro-relational communities of acquaintance and informal membership \citep{collins75, dimaggio87, purcell07, lizardo06a}.  I have argued that the deployment of cultural knowledge in concrete interaction is an important mechanism linking culture to status-based stratification.  In spite of its centrality, it has been a relatively under-studied phenomenon. This paper contributes to the study of culture and stratification by examining the socio-demographic distribution and cultural correlates of the joint tendency to talk to friends about culture consumption experiences.  

\paragraph{Mapping the social ecology of cultural talk.}
Empirically, I advance research on the role of culture-based talk as a constitutive component of lifestyle-based stratification in contemporary post-traditional societies.  I do this by examining the inter-relationship between aesthetic consumption, culture-talk and cultural taste. The results show that culture-based talk is a relatively common experience for the large majority of Americans, with almost four fifths reporting engaging in some form of culture-based talk using cultural knowledge from the five domains included in the NSS. The results also show however, (consistent with prior theory in cultural stratification) that culture-based talk is stratified by the same markers of social position (Age, Gender) and social status (Education) that previous research has shown also characterize audience segmentation in the corresponding cultural domains.  

In particular, I show that culture-talk in those domains characterized by large-scale commercial culture (Movies, Music) is strongly stratified by education (with the highly educated generating more culture-based talk) and in particular by age (with the young being more likely to use such culture in conversation). Culture-talk recruiting cultural knowledge from the higher-status ``non-profit'' arts is largely stratified by educational attainment, but only weakly differentiated across age-groups (net of education).  Literature is distinctive in that in addition to exhibiting the expected substantial rates of stratification by education it is also strongly stratified by gender (although all domains except for music show a female advantage in culture-based talk) as suggested by recent research on compositional differences of the ``reading class.''  

These results provide baseline expectations for the likely distribution of cultural content across socio-demographic groups.  For instance inter-generational segmentation in informal sociability is likely to be informally reproduced and generated by age-based segmentation on what type of culture is deployed in interaction \citep{milner04}, a mechanism that also serves to maintain divisions across educational groups \citep{collins75}.   By showing that culture-based talk is relatively commonplace but also socially structured in predictable ways, this paper takes a first step towards the empirical specification of an important but so far under-analyzed phenomenon in the sociology of taste.  

\paragraph{Heterogeneity in rates of culture-talk across taste groups.}
The results reported above show that a model that incorporates recent research and theory on ``intellectualization'' in cultural domains institutionalized as artworlds, generates a testable prediction: those persons more likely to be oriented toward high-status culture that will be more likely to generate culture-based talk. The results show that persons exhibiting high-status taste patterns in the domain of music, are more likely to generate culture-based talk within that domain.  This explains why we observe education-based stratification in the likelihood of generating culture-based talk in this domain in the absence of education-linked differences in the frequency of listening to music \citep{denora00}. These results contribute to our understanding of how the ``aesthetic disposition'' is mobilized in everyday interaction as a higher likelihood to engage in culture-based talk across domains. 

An important question for future research is whether a similar pattern can be observed in other domains associated with the commercial arts (e.g. Movies) or literature. The intellectualization hypothesis predicts that it is precisely those audiences more likely to gravitate towards the most ``intellectualized'' sub-genres (e.g.  in the case of movies, such genres as foreign, ``independent'' or documentary film) that are also more likely to use culture as an interactional resource.    Nevertheless, we must leave open the possibility that future studies that include a wider range of cultural practices and domains may lead to conclusions that heavily qualify the proposition connecting intellectualization and genre-legitimacy to culture-talk.  For instance within the realm of television-programming, it may be those genres endowed with the least amount of critical legitimacy (e.g. reality television) that generate the most culture-based talk, especially among the young.  The very limited research done on ``sports talk'' \citep{purcell07} suggests however, that ``intellectualization'' in the realm of sports (e.g. produced by adopting the critical vocabulary produced by sports journalists) does facilitate sports-based talk among educated sports enthusiasts, as suggested by the above analysis. 

\subsection{Contributions}
Empirically, this paper makes two main contributions.  First, I show that even if widespread, culture-talk is highly stratified by socio-demographic markers, in particular education (across both non-profit and commercial domains), age (primarily in commercial domains such as film and music) and gender (mainly when it comes to literature). second, I show that for domains in which status-based stratification based on talk is strong without corresponding levels of stratification based on consumption, culture-talk is stratified by membership in a given \emph{taste culture}.  Here, respondents more likely to belong to more artistically legitimated taste cultures engage in culture-talk more frequently.  Analytically, this research has important implications for our understanding of the link of culture and stratification in late-modern societies \citep{ollivier00}.  If the argument offered here is on the right track, then it becomes clear that for culturally-advantaged members of high-status groups the deployment of cultural capital \emph{in interaction} \citep{eliasoph_lichterman03} during face-to-face talk serves as an key \emph{relational mechanism} \citep{tilly98} that links social position to aesthetic consumption.  This mechanism explains both status-based advantages premised on the maintenance of (homo)social bonds, as argued in the classic Weberian tradition, as well as those social advantages tied to the ability to connect across distant social positions in the social structure \citep{erickson96, lizardo06a, lin99, pachuki_breiger10}.   Rethinking informal socialization as culture-mediated sociability allows us to reconsider the role of cultural content in the constitution of social relations \citep{ikegami05, mische03, mische10}.  

Analytically, this research has important implications for our understanding of the link of culture and stratification in late-modern societies \citep{ollivier00}.  If the argument offered here is on the right track, then it becomes clear that for culturally-advantaged members of high-status groups the deployment of cultural capital \emph{in interaction} \citep{eliasoph_lichterman03} during face-to-face talk serves as an key \emph{relational mechanism} \citep{tilly98} that links social position to aesthetic consumption.  This mechanism explains both status-based advantages premised on the maintenance of (homo)social bonds, as argued in the classic Weberian tradition, as well as those social advantages tied to the ability to connect across distant social positions in the social structure \citep{erickson96, lizardo06a, lin99, pachuki_breiger10}.   Rethinking informal socialization as culture-mediated sociability allows us to reconsider the role of cultural content in the constitution of social relations \citep{ikegami05, mische03, mische10}.  
